\begin{abstract}
动态操作要求机器人在场景持续变化的过程中，于合适的时刻和位姿完成接触。机器人应当记住和预测什么，才能满足这一要求？本文提出SE(3)-WAM，将带时间戳的历史、米制目标运动、场景预测和流匹配动作生成连接起来，并围绕三个问题开展受控研究：历史保留了哪些运动信息，未来预测采用什么变量与粒度，以及世界与动作分支如何交换信息。十四项动态任务从成功率、接触几何、时机和计算成本评价这些选择。诊断实验表明，历史更容易被读出、能够影响动作，并不必然改善控制；未来监督的收益也取决于其所在的表示与训练配置。通过匹配对照和预测干预，研究进一步区分辅助训练的作用与在线预判的效用。训练方案以真实机器人数据建立共同基座，分别初始化仿真和真机后训练；系统评测覆盖外部动态任务、静态兼容性与实际执行延迟。模型与评测协议共同将预测质量联系到动作真正可用的信息和时间。
\end{abstract}

\section{引言}

机器人抓取运动物体时，必须到达物体将要出现的位置，并满足接触所需的朝向与时机。向运动容器释放物体、穿过正在闭合的通道，也面临同样的要求。当前图像能够定位场景，却未必能区分其运动；计算与执行又使这幅图像逐渐过时。因此，动态操作既需要时间信息，也需要在实际生效时仍然合适的动作\cite{dynamicvla,domino,dynamicwam}。

\begin{figure}[ht]
\centering\includegraphics[width=\textwidth]{interaction_timing.pdf}
\caption{动态操作需要交互时刻的信息。两段历史到达相同的当前目标位姿(a)，却对应同一未来时刻的不同位姿(b)；成功接触还取决于朝向与周围间隙(c)。实线、虚线表示两种概念性运动，不是实测轨迹。}
\end{figure}

机器人应当记住和预测什么？重建所有可见变化是一种选择，但控制任务并不赋予每个像素相同的重要性。目标位姿决定接触几何，障碍物影响可达性与碰撞，遮挡物影响可见性；即使没有运动，静态几何也可能不可缺少。本文以交互需求组织信息：在足以支持动作的分辨率上保留运动与上下文。固定划分目标和环境，以及共享场景表示，是这一原则的两种可检验实现。

世界—动作模型联合学习未来场景与机器人动作\cite{dreamzero,openwam}，为上述问题提供了统一框架，但也混合了不同收益来源：额外监督、状态表示变化，以及行动前增加的计算。未来辅助损失可以在部署时不生成预测的情况下塑造特征\cite{fastwam}，生成的未来也可以成为动作的显式输入。要比较这些机制，需要区分策略看到了什么、学习预测什么，以及部署时实际使用什么。

SE(3)-WAM将带时间戳的历史编码为共同条件表示，由世界分支描述目标运动与场景演化，由动作分支在米制机器人坐标系中生成指令。三个问题构成设计研究：\textbf{Q1}研究历史保留的信息，\textbf{Q2}研究未来预测的范围、变量和分辨率，\textbf{Q3}研究参数共享、世界—动作访问关系与生成顺序。各项比较使用明确的输入和执行接口，并在选择完整模型前检验组件之间的相互影响。

诊断实验说明，这些区分会影响结论。结构化运动历史提高了参照策略的成功率，视觉历史更容易被读出却未必改善控制。未来监督同样依赖配方：稀疏轨迹在一种参照下没有明确总体收益，换用另一编码器与重建目标后，则出现集中于特定任务的较大收益。因此，研究需要考察完整的信息路径，不能仅依据历史或预测的数量评价设计。完整评测将这些诊断与匹配设计比较、真实机器人数据预训练，以及独立仿真和真机适配联系起来。

本文的贡献包括：

\begin{itemize}
\item 十四项动态操作任务及其评测协议，将任务成功与运动、接触和观测到动作的延迟联系起来。
\item 围绕历史、未来表示和世界—动作耦合的受控研究，通过干预区分信息保留、动作依赖与闭环效用。
\item 完整的SE(3)-WAM模型族与训练方案，结合结构化目标运动、可配置的场景预测，以及用于独立下游适配的共同机器人数据基座。
\end{itemize}

\section{相关工作}

\paragraph{动态操作与物理时间。}
动态抓取将运动预测、可达性与拦截联系起来\cite{akinola2021dynamic}。DOMINO/PUMA结合动态任务、历史光流与对象未来监督\cite{domino}；DynamicManip从静态示范扩展动态数据\cite{dynamicmanip}。DynamicVLA侧重连续推理与动作流式执行\cite{dynamicvla}，DynamicWAM将运动信息引入世界—动作学习\cite{dynamicwam}。时间特征预测\cite{dbctfp}、AHEAD的预测潜变量\cite{ahead}和ForeTime-VLA的未来token与事件时间监督\cite{foretime}，分别提供了预判交互的途径。RTC处理延迟下动作块的连续执行\cite{rtc}，Chronos Policy显式编码观测采集与动作应用的物理时间\cite{chronos}。这些工作提示，预测需要结合动作实际生效的时刻评价。

\paragraph{历史与对象信息。}
TraceVLA通过视觉轨迹提示表达时间信息\cite{tracevla}；HiF-VLA以运动表示组织历史与未来\cite{hifvla}。PACT-WAM为更早的观测分配更少空间tokens\cite{pactwam}，Long-WAM将长历史与因果视频预训练一并研究\cite{longwam}。历史编码的价值因而取决于时间支持和训练方式。PointWorld在不输入历史帧或速度时假设初始场景静止\cite{pointworld}；外部运动目标使初态信息成为Q1的关键因素。DOCIR区分操作对象、障碍物与机器人\cite{docir}，OA-WAM在世界—动作模型中保留可寻址的对象信息\cite{oawam}。本文先区分当前定位、特权运动与视觉历史，再比较编码方式，并将对象角色的组织方式作为实验因素。

\paragraph{未来的表示与作用。}
SE3-Nets用刚体变换描述场景运动\cite{se3nets}，SE3-Pose-Nets将结构化位姿动态用于控制\cite{se3posenets}。DSR-Net跨交互积累三维体素对象表示\cite{dsr}，PhysMani学习面向操作的三维世界表示\cite{physmani}。ATM从视频学习点轨迹\cite{atm}，Im2Flow2Act将生成的对象运动作为策略接口\cite{im2flow2act}，AffordTrajDP利用对象位姿传播接触可供性\cite{affordtraj}，FlowWAM在策略与世界模型两种模式中使用光流\cite{flowwam}。这些接口分别涉及期望运动、给定动作的后果或外部物体演化。DreamVLA与DreamWAM组合视觉、几何和语义未来目标\cite{dreamvla,dreamwam}，DC-WAM在RGB监督内部强调动态区域\cite{dcwam}，DELE-w0.5采用紧凑未来潜变量\cite{dele}。Q2在共同动作路径下比较相关目标的训练价值；其中的位姿输出沿用已有刚体运动参数化。

\paragraph{世界—动作交互与归因。}
预测通过不同接口进入控制。PointWorld预测候选动作条件下的三维点运动以评价动作\cite{pointworld}，NWM模拟并排序导航轨迹\cite{nwm}；CoT-VLA先生成视觉目标，再生成动作序列\cite{cotvla}，VPP以预测特征条件化动作学习\cite{vpp}。DreamZero联合建模世界与动作\cite{dreamzero}，Fast-WAM区分训练收益与部署时生成未来的作用\cite{fastwam}。本文参照D以未来监督训练，但动作不读取预测未来；Q3进一步改变访问关系和生成顺序。Wang等人的OpenWAM提供适配基线使用的Alpha基座\cite{openwam_alpha}，Yu等人的OpenWAM研究可组合的世界—动作程序\cite{openwam}。ACT、Diffusion Policy与$\pi_{0.5}$提供系统参照\cite{act,diffusionpolicy,pi05}，DROID提供真实机器人预训练轨迹\cite{droid}。系统比较和匹配设计对照分别回答不同问题。

\section{动态操作Benchmark}

\paragraph{物理时间中的交互。}
令$\mathcal H_t=\{(o_i,s_i,\tau_i)\}_{i\le t}$表示有限观测、机器人状态与时间戳历史。策略输出动作块$A_t=(a_{t,0},\ldots,a_{t,H_a-1})$。对于图像采集时刻$\tau_t^s$、观测到指令的延迟$\ell_t$和动作间隔$\Delta_a$，
\begin{equation}
\tau^{\mathrm{apply}}_{t,j}=\tau_t^s+\ell_t+j\Delta_a.
\end{equation}
预测时域以图像采集时刻为起点，接触则按动作实际应用时刻评价。增加未来动作token数量与预判物理运动是不同问题。

\paragraph{任务与示范。}
MuJoCo/robosuite套件\cite{mujoco,robosuite}包含十四项任务，覆盖移动目标与朝向敏感抓取、物理遮挡、运动容器放置、相位约束穿越、滚动或滑动物体，以及碰撞后适应。DB0 v0.6训练池包含11,200条示范，每任务800条。各任务按原生物理结果判定成功，并在总体得分中等权。附录A说明运动、事件、阈值与数据构造，补充图S1给出实际观察。

\paragraph{信息与执行条件。}
视觉部署条件使用因果传感器观测与机器人状态；需要目标运动的编码器使用在线估计量。当前位姿条件额外提供精确当前目标位姿，物体运动条件进一步提供区间运动学。这两种oracle条件用于区分表示能力与定位、跟踪误差。诊断筛选采用同步重规划，每次执行一条指令；T5的原生命令延迟为50\,ms，其余任务为零。部署评测让场景在实测推理等待期间继续演化，并单列人为延迟扫描，因此同步成绩与墙钟执行鲁棒性分别报告。

\paragraph{划分与指标。}
现有fixed140（每任务10场）及其1,260场补集共同构成开发池，各诊断表保留两部分身份。最终确认使用新的1,400场，每任务100场；场景种子独立，与训练、开发episode不重叠，任务生成器保持一致。模型选择只使用开发集，三个独立训练seed共享评测场景。对于二元结果$y_{mki}$，成功率为
\begin{equation}
\mathrm{SR}(m)=\frac1{14}\sum_{k=1}^{14}
\frac1{n_k}\sum_{i=1}^{n_k}y_{mki}.
\end{equation}
结果同时报告各seed成绩和配对差异；确认实验的区间分别重采样训练seed和场景，早期单seed筛选的区间只反映场景变化。接触时间、相对位姿、碰撞与失败类别、动作年龄和p95时延用于解释成功率。附录A规定其测量范围以及未发生接触的处理方式。

\section{SE(3)-WAM}

模型记为$(H,Z,B,G)$，分别表示历史编码器、未来表示、主干和世界—动作计算程序。未来观测提供训练监督；部署时，各配置仅以已有历史和生成状态为条件。具体参照采用八帧RGB与目标运动历史，以目标位姿和场景轨迹描述未来，使用共享注意力主干及辅助世界监督。Q1--Q3在保持相关输入和执行接口的前提下改变这些选择。

\begin{figure}[ht]
\centering\includegraphics[width=\textwidth]{model_design.pdf}
\caption{SE(3)-WAM：从带时间戳的历史到交互。因果编码器为世界与动作分支提供条件表示；世界分支描述米制目标位姿和场景运动，动作分支生成指令块。主干分组$B$与信息流$G$决定参数共享以及对生成世界状态的访问。未来观测提供训练监督；部署只使用查询时刻已有的观察和生成状态。下方分别展示辅助监督、联合生成和顺序生成，具体训练与采样定义见正文及附录E。}
\end{figure}

\subsection{带时间戳的历史与场景角色}

历史编码器输出$X_t=H(\mathcal H_t)$。每个观测token结合视觉内容、观测年龄$\tau_t^s-\tau_i$及可得的几何信息；机器人状态token包含末端位姿和实测夹爪开度；对象历史token保留米制位置、相对旋转、区间长度与可见性。因果估计只使用截至$t$的帧，离线平滑只用于监督，真实机器人数据同样遵循这一划分。

视觉状态参照使用冻结参数的SAM~2生成的掩码\cite{sam2}与FoundationPose RGB-D配准/跟踪\cite{foundationpose}。首个观察帧中的目标框初始化掩码，物体CAD或独立参考视图在episode开始前取得。跟踪按时间向前推进，不使用未来修正；无效估计产生显式缺失状态token与有效位，训练保持同一接口。附录F说明初始化、因果处理与共同输入条件。

比较包括当前RGB、堆叠RGB、codec运动向量、光流、压缩视频特征，以及按年龄或均匀分配的query压缩。它们以不同token成本提供时间信息。Oracle运动历史给出信息质量参照；视觉部署编码器仍需独立评价。Q1因此先区分当前定位与运动信息的可得性，再在共同时间支持下比较压缩方式。

根据任务指定当前目标，由此划分目标与环境角色。共享场景方案保留一个token bank，学习依赖任务的读出；显式角色方案加入目标/环境嵌入。各条件均保留静态环境观察，训练角色指定不向当前输入引入未来可见性。

\subsection{未来状态与预测分辨率}

世界表示包含目标运动和可配置的场景描述。给定当前目标位姿$(p_t,R_t)$及$\delta_h=0.05h$秒、$h=1,\ldots,8$，位姿标签为
\begin{equation}
d_h=p_{t+\delta_h}-p_t,\qquad
\phi_h=\Log(R_{t+\delta_h}R_t^\top)^\vee.
\end{equation}
重建使用$\hat p_{t+\delta_h}=p_t+\hat d_h$和
$\hat R_{t+\delta_h}=\Exp(\hat\phi_h^\wedge)R_t$。
确定性读出将16个meta tokens映射为累计平移与左相对旋转；有效标签上的损失为平移平方误差加0.1倍旋转测地距离平方。该表示采用米制几何，网络本身不预设等变性。

场景候选包括米制点位移、带可见性的稀疏图像轨迹和视频特征；目标候选还包括点轨迹或池化视觉特征。这些构成Q2的变量比较。另一个分配实验固定128条轨迹和八个未来时刻，分别使用64/64、96/32和32/96条目标/环境轨迹，总token数与角色信息相同。在相同支持上只施加目标或环境损失，用于检验预测范围。这样可以将空间分配效应与更换变量、增加token数量区分开。

在线生成状态$F_1$为标准化的未来场景位移（$196\times8\times3$）或轨迹坐标（$128\times8\times2$），统计量只来自训练数据。轨迹可见性由独立二分类头预测；GT有效掩码限制损失，不进入条件输入。目标位姿保持确定性读出。附录B、D定义几何支持、归一化与全部目标，使同一表示既可提供辅助监督，也可作为动作的显式输入。

\subsection{世界—动作训练与生成}

参照Transformer具有四层、宽768、八个注意力头。每项比较固定视觉前端：冻结的DINOv3编码器\cite{dinov3}，或组合实验中联合训练的CNN。D主干共享注意力投影、按流分离FFN；E主干将世界与动作分为两个参数组，各自具有注意力、调制和FFN。活动主干容量差控制在5\%内，前端与部署参数另行报告。

动作采用flow matching\cite{flowmatching}。对于归一化目标$a_1$、$a_0\sim\mathcal N(0,I)$和$u\sim\mathcal U[0,1]$，
\begin{equation}
\begin{aligned}
a_u&=(1-u)a_0+ua_1,\\
\mathcal L_A&=\mathbb E\|v_\theta^A(a_u,u,X_t,\mathcal C_G)-(a_1-a_0)\|^2.
\end{aligned}
\end{equation}
当动作不读取生成的世界状态时，$\mathcal C_G$为空；联合生成时，它包含带噪未来及其flow time，R5则包含已解码的第一阶段状态。世界生成状态采用相同插值与独立高斯源。总目标为
$\mathcal L=\mathcal L_A+\lambda_F\mathcal L_F+
\lambda_P\mathcal L_P+\lambda_V\mathcal L_V$，
只计算启用的目标及其有效支持；比较内固定权重和归一化。

Q3区分监督、访问与采样。R0仅训练动作；R1加入未来损失但不建立世界—动作访问边；R2允许动作读取世界状态，R3允许世界读取动作，R4双向读取。联合程序对每个生成模态使用五步流匹配。R5先生成世界再生成动作，R6反向执行，每阶段五步。顺序程序的第二阶段训练使用干净的示范第一阶段值，部署使用其生成值，因此比较包含teacher forcing造成的训练—部署差异。附录E给出全部注意力掩码与梯度路径。R1还对同一checkpoint比较仅动作与联合部署，测量生成未被动作读取的预测所需的成本。

每个动作块包含八条E7指令：基座系平移增量、左相对旋转向量及夹爪指令。控制器执行第一条后重新规划。条件缓存在单次查询内复用，后续查询刷新观测；参照动作求解器采用五步Euler。部署按物理时间戳执行，并将实测时延纳入评价。

\subsection{完整模型选择与机器人预训练}

模型选择面向完整配置。研究在保留的未来表示下回访历史编码，在不同信息流下回访未来表示，并在选定信息流下复核D/E。开发规则结合任务等权SR与p95时延，部署时延上限设为200\,ms，显存上限为共同设备容量的80\%。准确率与成本取舍采用2个百分点的实际差异界限及10\%的时延降低标准；未能明确区分的比较保留预先声明的参照（附录E）。新场景用于确认，不重新参与选择。

所得模型使用DROID轨迹预训练\cite{droid}。因果位姿与轨迹估计构造输入，离线几何估计提供带掩码的世界标签；标定、时间对齐、标签误差和监督覆盖率分别评价。实验在相同数据与训练曝光量下比较仅动作与世界—动作预训练，再改变容量和独立数据量。

同一个机器人数据基座分别初始化仿真与真机适配，两支使用独立训练数据、归一化和优化器，互不继承适配权重。无机器人预训练的对照保留相同的通用视觉初始化，以区分学习到的动态表示迁移与下游配方收益。

\section{实验}

评测覆盖完整系统、三个设计问题与机器人预训练迁移。表中横线表示缺失测量。诊断数字保留原开发划分和信息条件；附录A说明评分，B--H给出完整配置与结果。

\subsection{系统表现与比较协议}
\begin{table*}[t]
\centering\small
\caption{\tablelang{System evaluation with causal visual inputs. Dynamic14 reports task-macro success at zero and 150\,ms added delay; all success columns are percentages. External suites use their native success definitions. Latency is the end-to-end sensor-to-command p95 on the common deployment device. Each row uses the same task demonstrations and evaluation scenes; initialization is part of the system comparison.}{因果视觉输入下的系统评测。Dynamic14 报告无额外延迟及额外延迟 150\,ms 时的任务宏平均成功率；各成功率列单位均为百分比。外部评测使用各自的原生成功定义。时延为共同部署设备上从传感到命令输出的端到端 p95。各行使用相同任务示范和评测场景，初始化属于系统比较的一部分。}}
\label{tab:design_systems}
\begin{tabularx}{\textwidth}{@{}lXrrrrrr@{}}
\toprule
& & \multicolumn{2}{c}{Dynamic14} & \multicolumn{3}{c}{\tablelang{Transfer}{迁移}} & \tablelang{Latency}{时延} \\
\cmidrule(lr){3-4}\cmidrule(lr){5-7}\cmidrule(l){8-8}
\tablelang{System}{系统} & \tablelang{Initialization}{初始化} & \tablelang{0\,ms}{0\,ms} & \tablelang{150\,ms}{150\,ms} & DOMINO & LIBERO & \tablelang{Real}{实机} & p95 (ms) \\
\midrule
Diffusion Policy~\cite{diffusionpolicy} & \tablelang{Scratch}{从头训练} & \resultplaceholder & \resultplaceholder & \resultplaceholder & \resultplaceholder & \resultplaceholder & \resultplaceholder \\
$\pi_{0.5}$~\cite{pi05} & \tablelang{Published weights}{公开权重} & \resultplaceholder & \resultplaceholder & \resultplaceholder & \resultplaceholder & \resultplaceholder & \resultplaceholder \\
DreamVLA~\cite{dreamvla} & \tablelang{Published recipe}{公开配方} & \resultplaceholder & \resultplaceholder & \resultplaceholder & \resultplaceholder & \resultplaceholder & \resultplaceholder \\
\midrule
\ours & \tablelang{Scratch}{从头训练} & \resultplaceholder & \resultplaceholder & \resultplaceholder & \resultplaceholder & \resultplaceholder & \resultplaceholder \\
\ours & \tablelang{DROID pretraining}{DROID 预训练} & \resultplaceholder & \resultplaceholder & \resultplaceholder & \resultplaceholder & \resultplaceholder & \resultplaceholder \\
\bottomrule
\end{tabularx}
\end{table*}


系统主表比较Dynamic14、延迟鲁棒性、外部基准和真实操作。最终Dynamic14使用新的1,400场确认集及视觉部署条件。各基线保留公开模型组织方式，采用共同观察、示范与物理动作语义，并获得相同下游调参机会；预训练、训练参数、累计样本量和端到端成本另列。SE(3)-WAM有无机器人预训练的对照还共享下游数据与曝光量。

既有适配系统结果保留在附录B：DreamVLA独立二十轮和四十轮配方在1,260场开发补集上分别为79.37\%和79.60\%，所选特权信息组合为66.59\%。这些成绩对应不同输入和预算，不迁入新的视觉测试表。

\subsection{Q1：什么历史有助于控制？}
% Source: Research/HISTORY1/HISTORY1.md, read 2026-10-16.
\begin{table}[t]
\centering\small
\caption{Visual history given current GT target pose and robot history.
Historical target kinematics are removed. SR uses the 1,400-scene development
total; $\Delta$ is relative to current image, in percentage points. Appendix C
reports partition counts and four unfinished conditions separately.}
\label{tab:q1_results}
\begin{tabular}{@{}lrr@{}}
\toprule
Visual input & SR /1,400 (\%) & $\Delta$ (pp) \\
\midrule
Current image & 53.43 & 0.00 \\
Window & 47.29 & $-6.14$ \\
Codec-MV & 52.36 & $-1.07$ \\
Flow & 51.43 & $-2.00$ \\
PACT-Age & 45.29 & $-8.14$ \\
\bottomrule
\end{tabular}
\end{table}


Q1固定当前目标位姿与机器人历史，比较八种视觉历史编码，并加入重新训练的GT运动参照。当前图像达到53.43\%，七种视觉历史编码均未超过其点估计；GT运动达到58.86\%，提高5.43个百分点，但相位穿越少成功83场。该协议研究已知定位时历史的增量价值。

三项测量从不同层面解释结果。冻结读出评价保留的信息：Window对350\,ms前位置的误差由10.89降至8.46\,mm。动作干预评价输入使用：重复或反转过去图像会改变七种历史模型的动作，而两个当前图像参照保持不变。闭环成功率评价控制效用：Window的读出更准确，SR却低6.14个百分点。因此，在这一配方下，信息保留和动作敏感性并不等于更好的控制。

固定GT运动端点，将过去位姿替换为当前位姿并重新计算运动学，成功数由824降至9/1,400，即下降58.21个百分点（95\%配对区间56.50--59.86）。210个探针窗口的首条动作全部改变，平移变化均值为18.40\,mm；恒等干预保持动作一致，并复现140场控制评测的全部成功/失败标签。这是推理时历史失配的效应；独立训练的当前图像参照仍达到53.43\%。

匹配确认将采样间隔、遮挡与视觉/估计状态历史交叉比较。重复历史保持token数量，反序保留外观但改变运动顺序，年龄扰动检验时间信息。闭环干预、接触误差、编码耗时与运动读出共同报告（附录C），用于区分缺少信息、编码不准和未能有效使用运动信息。

\subsection{Q2：应当预测什么未来？}

\begin{figure}[ht]
\centering\includegraphics[width=\textwidth]{future_supervision.pdf}
\caption{1,260个配对开发场景上的未来监督比较：十四任务各90场，一个训练seed。(a)十五组相对P00（62.78\%）的SR差值与任务内配对bootstrap的逐项95\%区间；(b)所列比较新增与损失的成功场景。P05经27项比较校正后表现不利，P07/P09未建立明确总体收益。fixed140与逐任务结果见附录D。}
\end{figure}

\paragraph{预测变量。}
十五组筛选固定当前RGB、DINO-L、目标运动学、训练曝光量与仅动作部署。四轮训练包含40,104次更新、641,664个样本块，全局批量16。无未来P00成功791/1,260场，场景流、目标位姿及其组合分别为782、789、786；场景轨迹加位姿P07与场景流加目标轨迹P09分别为807、808。两者配对新增/损失成功为85/69和91/74，均未显示明确总体优势。稠密二维位移加位姿P05为746/1,260，比P00低3.57个百分点，校正$p=0.00844$。该不利结果属于具体表示配置；较有竞争力的稀疏轨迹同样使用二维坐标。

\paragraph{未来表示与视觉编码器的相互影响。}
十二组比较交叉三类未来目标、DINO-S/CNN与当前目标重建，均保留完整视觉窗口和特权运动。场景轨迹加位姿、CNN与重建的候选I获得935/1,400（66.79\%）；匹配的无未来与场景流/位姿控制分别为810、861，差值为8.93和5.29个百分点。相对后者净增的74场中，72场来自相位穿越。I由这一开发池选出，其第二个训练seed达到61.21\%。附录D给出全部组合；确认比较对I和两个控制使用共同训练seed及未参与选择的场景。

\paragraph{范围与粒度。}
十五种目标同时改变了变量和支持；同族轨迹实验则单独比较目标、环境、均匀目标/环境、目标细与环境细。各条件共享当前目标身份、时间支持和128个tokens。共享bank对照保持采样点不变，移除显式角色嵌入，用于检验固定角色划分是否优于任务条件读出。位姿精度、可见轨迹误差、接触误差与SR一起报告（附录D）。几何误差下降而任务得分不变，意味着预测改善，不能单独作为控制表示更优的证据。

\subsection{Q3：世界与动作应当怎样交互？}
\begin{table}[t]
\centering\small
\caption{\tablelang{Backbone and information flow. The upper contrast varies parameter grouping at fixed access. The lower block fixes the backbone $B^\star$ selected by the development rule and varies supervision, access, or generation order, with $B^\star$/R1 as its reference. SR is the three-seed mean task-macro success (\%). End-to-end p95 includes the generated branches used by each program.}{主干与信息流。上部在固定访问关系下改变参数分组。下部固定按开发集规则选出的主干 $B^\star$，以 $B^\star$/R1 为参照，改变监督、访问边或生成顺序。SR 为三个种子的任务宏平均成功率均值（\%）。端到端 p95 包含各生成程序实际使用的分支。}}
\label{tab:design_study}
\begin{tabularx}{\linewidth}{@{}Xrr@{}}
\toprule
\tablelang{Configuration}{配置} & SR $\uparrow$ & p95 (ms) $\downarrow$ \\
\midrule
D/R1: \tablelang{shared attention, split FFNs}{共享注意力，分支 FFN} & \resultplaceholder & \resultplaceholder \\
E/R1: \tablelang{world/action grouping}{世界/动作参数分组} & \resultplaceholder & \resultplaceholder \\
\midrule
$B^\star$/R0: \tablelang{action loss only}{仅动作损失} & \resultplaceholder & \resultplaceholder \\
$B^\star$/R1: \tablelang{auxiliary, action only}{辅助监督，仅生成动作} & \resultplaceholder & \resultplaceholder \\
$B^\star$/R2: \tablelang{joint, $F\!\to\! A$}{联合，$F\!\to\! A$} & \resultplaceholder & \resultplaceholder \\
$B^\star$/R4: \tablelang{joint, $F\!\leftrightarrow\! A$}{联合，$F\!\leftrightarrow\! A$} & \resultplaceholder & \resultplaceholder \\
$B^\star$/R5: \tablelang{future, then action}{先未来，后动作} & \resultplaceholder & \resultplaceholder \\
\bottomrule
\end{tabularx}
\end{table}


Q3首先固定历史、未来目标与仅生成动作的部署方式，在辅助世界监督下比较D/E；再固定主干，比较R0--R6（完整矩阵见附录E）。R0/R1检验世界监督的训练作用，R1/R2检验动作读取生成世界状态的价值，R2/R4检验世界读取动作的作用，R2/R5检验顺序生成的准确率与成本取舍。五步联合与五加五步顺序生成匹配每模态步数，总计算量仍需实测。

对读取世界状态的模型，冻结策略后将预测替换为匹配的零运动、打乱或有效支持上的oracle未来。接触前外生运动可直接采用oracle；接触后由同一仿真状态和候选动作生成对应未来。测量包括动作位移、接触时机与配对SR。R1的仅动作/联合部署则测量生成未被读取预测的成本。这些对照区分训练作用、预测依赖与有益的在线预判。

\subsection{预训练、迁移与部署}

\paragraph{机器人数据预训练。}
实验比较无机器人预训练、仅动作预训练与世界—动作预训练，同时采用共同下游配方及等预算调参比较。主干容量约为参照的$0.5,1,2,4$倍；嵌套数据比例为$1/8,1/4,1/2,1$。四条完整数据曲线，加上两个容量各三个较小子集，构成十条主曲线。独立曝光量比较区分新经验与重复训练；10/25/50/100\%下游示范比例检验数据效率。附录F给出端点、标签质量对照、学习曲线与结果矩阵。

\paragraph{外部迁移。}
单臂DOMINO评价外部动态操作，LIBERO\cite{libero}评价静态多任务控制兼容性。真机覆盖运动抓取、运动容器放置、限时通道与碰撞后重新拦截；每任务交叉三档速度、三档延迟，并使用配对初始条件。每个下游分支独立从同一机器人数据基座初始化。域内归一化、动作接口适配与原生任务评分一并报告（附录F）。

\paragraph{延迟与接触。}
评测包含0/50/100/150/200\,ms人为延迟和实测时延执行模式。时间戳覆盖采集、预处理、模型推理、指令发送与应用；预测和动作时域始终以物理时间表示。既有未来筛选的模型内调用均值为177.071--184.212\,ms，p95为178.726--185.133\,ms，不含传感器与命令传输；系统主表使用端到端时延。按延迟报告SR和接触误差，可以检验预测的计算成本是否抵消其预判价值。

\paragraph{失败分析。}
十五种未来条件均未成功完成刚体运动容器放置，完整约束穿越也接近零。监督审计发现，800条完整穿越示范的释放只出现在动作块末位，而评测执行首位。附录将监督覆盖与表示失败分别分析，并对齐定性预测、实际指令和物理结果。这些诊断指出了仅改善场景预测难以解决的控制问题。

\section{结论}

SE(3)-WAM围绕交互时刻所需的信息研究动态操作，将带时间戳的历史、结构化目标运动、场景预测与动作生成纳入明确的训练和部署接口。Benchmark与实验区分信息质量、辅助监督、在线预测使用和计算成本。诊断证据说明了这些区分的必要性：可读出的历史未必改善控制，未来监督的收益也可能依赖周围的表示配置。完整评测通过匹配确认、机器人数据预训练与独立下游适配检验这些设计。
